% !TEX program = xelatex
% ============================================================
%  خطاب توصية — Recommendation Letter Template
%  قالب عربي (RTL) لخطاب توصية أكاديمي أو مهني.
%  Compile with: xelatex main.tex (run twice)
% ============================================================
%  Features:
%    - Right-to-left (RTL) layout with Amiri font
%    - Institution letterhead placeholder
%    - Addressed to admissions/hiring committee
%    - Relationship-to-candidate opening paragraph
%    - Three quality paragraphs with example placeholders
%    - Strong closing recommendation and signature block
% ============================================================
\documentclass[12pt,a4paper]{article}

% --- Page geometry ---
\usepackage[a4paper,margin=2.5cm,top=2.5cm,bottom=2.5cm,headheight=16pt]{geometry}

% --- Standard packages (load ALL before bidi) ---
\usepackage{graphicx}
\usepackage{array}
\usepackage{booktabs}
\usepackage{tabularx}
\usepackage{enumitem}
\usepackage{titlesec}
\usepackage{fancyhdr}
\usepackage{setspace}
\usepackage{xcolor}
\usepackage{parskip}

% --- BIDI (must come before polyglossia, after standard packages) ---
\usepackage{bidi}

% --- Font specification (requires bidi) ---
\usepackage[no-math]{fontspec}

% --- Polyglossia (language support) ---
\usepackage[quiet]{polyglossia}

% --- Fonts ---
\setmainfont[Scale=1]{Amiri-Regular.ttf}
\newfontfamily\arabicfont[Script=Arabic,Scale=0.95]{Amiri-Regular.ttf}
\newfontfamily\englishfont[Scale=0.95]{TeX Gyre Termes}

% --- Languages ---
\setdefaultlanguage[calendar=gregorian,hijricorrection=1]{arabic}
\setotherlanguage[variant=british]{english}

% --- Hyperref (ALWAYS last) ---
\usepackage[breaklinks,hidelinks]{hyperref}
\hypersetup{pdftitle={خطاب توصية},pdfauthor={اسم المرسل}}

% --- Line spacing ---
\setstretch{1.4}
\setlength{\parskip}{0.6em}
\setlength{\parindent}{0pt}

% --- Header/Footer ---
\pagestyle{fancy}
\fancyhf{}
\fancyhead[R]{\small\itshape خطاب توصية}
\fancyhead[L]{\small اسم المؤسسة}
\fancyfoot[C]{\small\thepage}
\renewcommand{\headrulewidth}{0.3pt}

% ============================================================
%  Document
% ============================================================
\begin{document}

% --- Letterhead placeholder (institution) ---
\begin{center}
{\LARGE\bfseries اسم المؤسسة}\\[0.3em]
% TODO: استبدل باسم المؤسسة/الجامعة الكامل
{\large اسم الجامعة / المؤسسة}\\[0.2em]
% TODO: استبدل بعنوان المؤسسة
{\small العنوان: المدينة، الحي، الشارع، الرمز البريدي}\\[0.2em]
% TODO: استبدل ببيانات التواصل
{\small الهاتف: \LR{+000-000-000-0000} \,\textbar\, البريد: \LR{info@institution.example}}
\end{center}

\vspace{0.5em}
\hrule
\vspace{1.5em}

% --- Date ---
\begin{flushright}
% TODO: حدّث التاريخ
\today
\end{flushright}

\vspace{1em}

% --- Recipient block ---
\begin{flushright}
% TODO: حدّد الجهة المُرسَل إليها (لجنة القبول أو إدارة التوظيف)
إلى لجنة القبول / إدارة التوظيف المحترمة،\\[0.2em]
% TODO: اسم الجهة المستلمة
\textbf{اسم الجهة / البرنامج}\\[0.2em]
% TODO: عنوان الجهة المستلمة
العنوان: المدينة، الحي، الشارع
\end{flushright}

\vspace{1em}

% --- Subject line ---
\begin{flushright}
\textbf{الموضوع: خطاب توصية للسيد/ة \textit{اسم المرشح/ة}}
\end{flushright}

\vspace{0.5em}

% --- Salutation ---
\begin{flushright}
المحترمون،
\end{flushright}

\vspace{0.5em}

% --- Body: Paragraph 1 — Relationship to candidate ---
% TODO: اذكر طبيعة العلاقة بالمرشح ومدتها
تحية طيبة وبعد،

يسعدني أن أكتب إليكم هذا الخطاب لتوصية السيد/ة \textbf{اسم المرشح/ة} للقبول في \textbf{البرنامج / الوظيفة} لدى \textbf{اسم الجهة}. عرفتُ المرشح/ة بصفتي \textbf{صفة المُوصي: أستاذ / مشرف / مدير مباشر} على مدى \textbf{مدة العلاقة بالسنوات} سنوات، خلالها تمكّنت من ملاحظة قدراته الأكاديمية والمهنية عن قرب.

% --- Body: Paragraph 2 — Quality 1 (academic/professional competence) ---
% TODO: اذكر ميزة محددة مع مثال قابل للقياس
يتميّز السيد/ة \textbf{اسم المرشح/ة} بتفوّقٍ واضح في \textbf{المجال/المادة}. فعلى سبيل المثال، خلال \textbf{المشروع / المقرر / المهمة}، أظهر قدرةً عالية على \textbf{المهارة المحددة}، حيث \textbf{إنجاز محدد قابل للقياس: حقق نتيجة / قاد فريق / حلّ مشكلة}. وقد عكست هذه التجربة تمكّنه من \textbf{مهارة قابلة للقياس} وفهماً عميقاً لأساسيات \textbf{المجال}.

% --- Body: Paragraph 3 — Quality 2 (personal/interpersonal qualities) ---
% TODO: اذكر ميزة شخصية أو قيادية مع مثال
إلى جانب الكفاءة العلمية، يتميّز بصفاتٍ شخصية رفيعة. فهو/هي يتمتّع بـ\textbf{الصفة: النزاهة / روح المبادرة / العمل الجماعي}، وقد تجلّى ذلك في \textbf{موقف محدد: قيادة فريق / تنظيم فعالية / حلّ نزاع}. كما يُظهر التزاماً عالياً بالمواعيد واحتراماً للزملاء، مما يجعله/ها عنصراً إيجابياً في أي بيئة عمل أو دراسة.

% --- Body: Paragraph 4 — Quality 3 (growth / impact) ---
% TODO: اذكر نمو أو أثر إضافي مع مثال
ومما يستحق الذكر أيضاً، قدرته/ها على التعلّم السريع والتكيّف مع المتغيرات. خلال \textbf{فترة أو حدث}، تطوّر من \textbf{مستوى ابتدائي} إلى \textbf{مستوى متقدّم} في \textbf{مهارة أو مجال}، وأحدث أثراً ملموساً على \textbf{الفريق / القسم / المشروع}. هذه القدرة على النمو المستمر تُعدّ مؤشراً قوياً على نجاحه/ها في المرحلة القادمة.

% --- Closing: strong recommendation ---
\vspace{0.5em}
% TODO: عزّز التوصية بدرجة التأييد المناسبة
بناءً على ما تقدّم، أوصي بقبول السيد/ة \textbf{اسم المرشح/ة} بقوةٍ في \textbf{البرنامج / الوظيفة}، وأرى أنه/ها سيكون/ستكون إضافة قيّمة لبرنامجكم. وأنا على استعدادٍ لتقديم أي معلومات إضافية عند الطلب.

% --- Closing phrase ---
\vspace{0.5em}
\begin{flushright}
وتفضلوا بقبول فائق الاحترام والتقدير،
\end{flushright}

\vspace{2em}

% --- Signature block ---
\begin{flushright}
% TODO: التوقيع اليدوي أعلى الاسم
\rule{5cm}{0.4pt}\\[0.3em]
\textbf{اسم المُوصي}\\[0.2em]
% TODO: المسمى الوظيفي للمُوصي
المسمى الوظيفي / اللقب الأكاديمي\\[0.2em]
% TODO: اسم المؤسسة
اسم المؤسسة / الجامعة\\[0.2em]
% TODO: بيانات التواصل
\LR{+000-000-000-0000} \,\textbar\, \LR{recommender@institution.example}
\end{flushright}

\end{document}
